\documentclass[11pt]{article}
\usepackage{preamble}

\newcommand{\IconRect}{[RECT]}
\newcommand{\IconCheck}{[CHECK]}
\newcommand{\IconMic}{[MIC]}
\newcommand{\IconClipboard}{[CLIPBOARD]}

\newcommand{\phasetag}[1]{%
  {\small\itshape Tag: #1\par}
}

\newcommand{\phaseitems}[1]{%
  \begin{itemize}[leftmargin=1.5em,itemsep=2pt,topsep=2pt]
  #1
  \end{itemize}%
}

\newcommand{\teacheranswerbox}[2]{%
  \begin{callout}{#1}{calloutgreen}
  #2
  \end{callout}
}

\begin{document}

\packetheading{Lesson 4-1 \textperiodcentered\ Period 1 \textperiodcentered\ Teacher Edition}{Inverse Variation --- Introduction}

\sectionbanner{OBJECTIVES}
{\small
\renewcommand{\arraystretch}{1.25}
\noindent\begin{tabular}{|p{1.8in}|p{4.8in}|}
\hline
\textbf{Math Objective} & Identify when two variables vary inversely by testing whether their product \(xy\) is constant, then write the inverse-variation equation \(y = k/x\) and use it to predict unknown values. \\ \hline
\textbf{Language Objective} & Students will use the frame ``As \_\_\_ increases, \_\_\_ decreases; their product is \_\_\_, so the equation is \(y = \_\_\_/x\).'' to explain an inverse-variation relationship. \\ \hline
\textbf{Essential Understanding} & Two variables vary inversely when their product is a constant \(k\). The equation \(y = k/x\) is the algebraic form; \(xy = k\) is the structural form. Graphically, the relationship is a hyperbola (the reciprocal function). \\ \hline
\end{tabular}
}

\sectionbanner{TOPIC GOALS / ESSENTIAL QUESTION / MATERIALS}
{\small
\renewcommand{\arraystretch}{1.25}
\noindent\begin{tabular}{|p{1.8in}|p{4.8in}|}
\hline
\textbf{Topic Goals} & Students will distinguish inverse variation from direct variation using product vs ratio tests, and write the equation \(y = k/x\) from two given values. All student-facing items are Savvas-sourced (bank IDs in teacher edition). \\ \hline
\textbf{Essential Question} & How are inverse variations related to the reciprocal function? \\ \hline
\textbf{Materials} & Pencils; student packet; calculator (optional). No Chromebooks required for Period 1. \\ \hline
\end{tabular}
}

\begin{callout}{\IconPin\ PRE-FLIGHT NOTE (read before class)}{calloutpink}
Explore is 7 bank items: Try-It 1 \(\rightarrow\) \#9 \(\rightarrow\) \#14 \(\rightarrow\) Try-It 2 \(\rightarrow\) \#13 \(\rightarrow\) \#16 \(\rightarrow\) \#10 (extension). Paced for \(\sim\)4--5 min per item.\\
If pairs finish early, push them to the Optional Reinforcement: Model \& Discuss Parts B + C (verbal/TPS), then Practice \#15.\\
If pace slows, cut \#10 (Construct Arguments) --- it's the DOK-2 stretch; the core DOK-1 arc (Try-It 1 \(\rightarrow\) \#9/\#14 \(\rightarrow\) Try-It 2 \(\rightarrow\) \#13/\#16) is the minimum viable set.
\end{callout}

\phasetag{notice \& wonder}
\frameworkphaseheader{Do Now}{2}{5}{%
\phaseitems{
  \item Project the two rectangles (or use printed Do Now sheet).
  \item Circulate silently. Do NOT teach during this phase.
  \item Collect at 5 min sharp.
}%
}{%
\phaseitems{
  \item Write 1 notice + 1 wonder.
  \item Sketch a third rectangle with area 144.
}%
}{%
\phaseitems{
  \item (None during the phase --- teacher harvests answers in Launch.)
}%
}{Solo teacher: circulate silently. Resist the urge to give hints.}

\begin{callout}{\IconRect\ TWO RECTANGLES --- Both have area = 144 square units}{calloutyellow}
Rectangle 1: length = 72, width = 2\\
Rectangle 2: length = 24, width = 6\\

A) Notice: Write something you notice about these two rectangles.\\

B) Wonder: Write something you wonder about rectangles with area 144.\\

C) Sketch: Draw a third rectangle with area 144. Label its length and width.
\end{callout}

\bankitem{Do Now (4-1-savvas-model-discuss-lesson-4-1-launch)}{%
The two rectangles shown both have an area of 144 square units (\(72 \times 2\) and \(24 \times 6\)). (A) Sketch as many other rectangles as you can that have the same area. Organize and record your data for the lengths and widths of the rectangles. (B) Use Structure: Considering rectangles with an area of 144 square units, what happens to the width of the rectangle as the length increases? (C) Examine at least five other pairs of rectangles, each pair sharing the same area. How would you describe the relationship between the lengths and widths?
}{}

\begin{callout}{\IconCheck\ ANTICIPATED STUDENT RESPONSES}{calloutgreen}
NOTICE (typical): ``Both rectangles have the same area.'' ``One is long and thin, the other is shorter and wider.'' ``When the length is bigger, the width is smaller.''\\
WONDER (typical): ``Can any rectangle have area 144?'' ``What's the rule?'' ``Is there a formula?''\\
SKETCHES: students might draw \(1 \times 144\), \(4 \times 36\), \(8 \times 18\), \(12 \times 12\), \(3 \times 48\), etc. All valid.\\
KEY MOVE in Launch: surface the insight that length \(\times\) width = 144 for every rectangle --- this is \(xy = k\).
\end{callout}

\sectionbanner{LAUNCH --- Inverse Variation (teacher models on screen)}

\begin{callout}{\IconBook\ INVERSE VARIATION RULES (keep printed on your page)}{calloutgreen}
\textbullet\ Two variables vary INVERSELY when their product is a constant \(k\):\\
\hspace*{1.6em}\(xy = k\) \hspace{1.4em} or equivalently \hspace{1.4em} \(y = k/x\)\\
\textbullet\ The constant \(k\) is the CONSTANT OF VARIATION.\\
\textbullet\ As one variable INCREASES, the other DECREASES proportionally.\\
\textbullet\ Example: if \(x = 10\) and \(y = 3\) vary inversely, then \(k = xy = 30\).\\
\hspace*{1.6em}Equation: \(y = 30/x\). When \(x = -6\): \(y = 30/(-6) = -5\).
\end{callout}

\begin{callout}{\IconWarn\ COMMON ERROR}{calloutred}
Do NOT write \(y/x = k\) --- that is DIRECT variation.\\
Inverse variation multiplies, direct variation divides.\\
Test: compute \(xy\) for several points. If constant \(\rightarrow\) inverse.
\end{callout}

\phasetag{teacher models Example 1 + Example 2}
\frameworkphaseheader{Launch}{1--2}{12}{%
\phaseitems{
  \item Bridge from Do Now: ``You all noticed length \(\times\) width = 144. That's a RULE, and it has a name: inverse variation.''
  \item Project Example 1 (Savvas Topic 4-1). Think aloud: for each table, compute \(xy\) for each column. Is it constant? \(\rightarrow\) YES = inverse; NO = not.
  \item Project Example 2. Think aloud: given \(x = 10\), \(y = 3\), find \(k = xy = 30\). Write \(y = 30/x\). Substitute \(x = -6\). Answer: \(y = -5\).
  \item Pause 1\(\times\) for 30-sec partner-check: ``What's \(k\) for your Do Now rectangles?'' (Answer: \(k = 144\).)
  \item Do NOT solve Try-Its. Release promptly at minute 17.
}%
}{%
\phaseitems{
  \item Watch silently through Example 1.
  \item During partner-check: ``What's \(k\) for area = 144?'' (144)
  \item Copy the Inverse Variation Rules card into their page margin.
  \item Note the COMMON ERROR warning: \(y/x = k\) is DIRECT variation.
}%
}{%
\phaseitems{
  \item ``What's \(xy\) for column 1? Column 2? Column 3? Is it constant?''
  \item ``If two variables have \(xy = 30\) always, what is \(k\)?''
  \item ``How do I write \(y\) in terms of \(x\) if \(xy = 30\)?''
}%
}{Project, think aloud. Do NOT give students Try-Its to work until minute 17.}

\begin{callout}{\IconMic\ TEACHER SCRIPT}{calloutblue}
1. ``From the Do Now, you noticed length \(\times\) width = 144 for every rectangle. That's inverse variation --- a pattern.''\\
2. Example 1 walk-through: ``Watch me. For each column, I compute \(xy\). If it's the same number for every column, the two variables vary inversely.''\\
3. Example 2 walk-through: ``Now I'll write the EQUATION. \(k = xy\), so if \(x = 10\) and \(y = 3\), then \(k = 30\). The equation is \(y = k/x\), so \(y = 30/x\). When \(x = -6\), \(y = 30/(-6) = -5\).''\\
4. Common-error flag: ``Don't mix up \(y/x = k\) (direct) with \(xy = k\) (inverse). Inverse multiplies.''\\
5. Release: ``Now you try. Try-It 1 asks if two tables are inverse variation. Try-It 2 asks you to write the equation.''
\end{callout}

\sectionbanner{EXPLORE --- Think-Pair-Share (33 min)}
\textit{Work with a partner. Use the Inverse Variation Rules above for every problem. Move through items in order --- each builds on the last.}

\phasetag{Think-Pair-Share + Practice}
\frameworkphaseheader{Explore}{1--2}{33}{%
\phaseitems{
  \item 5 laps circulation across 33 min (\(\sim\)6--7 min each).
  \item Lap 1: Try-It 1 --- press pairs to compute \(xy\) for every column.
  \item Lap 2: \#9 + \#14 --- check that pairs notice NOT every table is inverse.
  \item Lap 3: Try-It 2 --- press pairs to state \(k\), then write the equation.
  \item Lap 4: \#13 + \#16 --- quick check on generalization and application.
  \item Lap 5: \#10 (Construct Arguments) --- stretch. Cut if pace slows.
  \item Do not give answers. Use sentence frame to press justification.
}%
}{%
\phaseitems{
  \item Pairs move in order through 7 items: 2 Try-Its + 5 Practice.
  \item For each table (\#14, \#15 later): compute \(xy\) for EVERY column before answering.
  \item For Practice \#10 (why can \(x=0\) not be in the domain?): justify with the equation \(y = k/x\).
  \item Justify with sentence frame before writing.
}%
}{%
\phaseitems{
  \item What's \(xy\) for each column? Is it constant?
  \item If \(xy\) is NOT constant, what kind of variation is it? (Check: is \(y/x\) constant = direct variation.)
  \item What's the value of \(k\)? How do you know?
  \item Can you write the equation \(y = k/x\)?
  \item For \#10: if \(y = k/x\) and \(x = 0\), what happens? (Division by zero = undefined.)
}%
}{Solo teacher: 5 laps. Press pairs to justify using the rule card. No answers.}

\textbf{Bank-item labels (in order):}
\begin{enumerate}[leftmargin=1.6em,itemsep=2pt,topsep=2pt]
  \item Try It 1 --- Identify inverse variation
  \item Practice \#9 --- Generalize
  \item Practice \#14 --- Another table check
  \item Try It 2 --- Write the equation
  \item Practice \#13 --- Generalize: write \(k\) in terms of \(x\) and \(y\)
  \item Practice \#16 --- Apply the equation
  \item Practice \#10 --- Construct Arguments (extension)
\end{enumerate}

\bankitem{Try It 1 --- Identify inverse variation}{%
Try It 1. Determine if each table of values represents an inverse variation.\\[0.35em]
(a) \(x: 1,2,3,5,6,15;\ y: 25.5, 12.75, 8.50, 5.10, 4.25, 1.70.\)
\begin{center}
\renewcommand{\arraystretch}{1.15}
\begin{tabular}{|c|c|c|c|c|c|c|}
\hline
\(x\) & 1 & 2 & 3 & 5 & 6 & 15 \\ \hline
\(y\) & 25.5 & 12.75 & 8.50 & 5.10 & 4.25 & 1.70 \\ \hline
\end{tabular}
\end{center}
(b) \(x: 6.6, 5.5, 4.4, 3.3, 2.2, 1.1;\ y: 3, 5, 7, 9, 11, 13.\)
\begin{center}
\renewcommand{\arraystretch}{1.15}
\begin{tabular}{|c|c|c|c|c|c|c|}
\hline
\(x\) & 6.6 & 5.5 & 4.4 & 3.3 & 2.2 & 1.1 \\ \hline
\(y\) & 3 & 5 & 7 & 9 & 11 & 13 \\ \hline
\end{tabular}
\end{center}
}{}

\teacheranswerbox{\IconCheck\ ANSWER / ANTICIPATED TALK}{%
Answer key: (a) all products = 25.5 \(\rightarrow\) yes, inverse variation; (b) products not constant (\(6.6\cdot 3 = 19.8\); \(5.5\cdot 5 = 27.5\); ...) \(\rightarrow\) no. Use during Explore as check-for-understanding after Example 1.
}

\bankitem{Practice \#9 --- Generalize}{%
9. Generalize: Just from looking at the table of values, how can you determine that the data do NOT represent an inverse variation? Table - \(x: -2, 2, 4, 6, 8, 10;\ y: -6, 6, 12, 18, 24, 30.\)
\begin{center}
\renewcommand{\arraystretch}{1.15}
\begin{tabular}{|c|c|c|c|c|c|c|}
\hline
\(x\) & -2 & 2 & 4 & 6 & 8 & 10 \\ \hline
\(y\) & -6 & 6 & 12 & 18 & 24 & 30 \\ \hline
\end{tabular}
\end{center}
}{}

\teacheranswerbox{\IconCheck\ ANSWER / ANTICIPATED TALK}{%
Answer key (from TE): As the \(x\)-values increase, the \(y\)-values also increase. (Additionally, this IS a direct variation with \(k = 3\) since \(y/x = 3\) for every pair.)
}

\bankitem{Practice \#14 --- Another table check}{%
14. Do the tables of values represent inverse variations? Explain. See Example 1. Table - \(x: -1/4, -1/2, 1/3, 2, 5, 11;\ y: -9/2, -9, 6, 36, 90, 198.\)
\begin{center}
\renewcommand{\arraystretch}{1.15}
\begin{tabular}{|c|c|c|c|c|c|c|}
\hline
\(x\) & \(-\frac14\) & \(-\frac12\) & \(\frac13\) & 2 & 5 & 11 \\ \hline
\(y\) & \(-\frac92\) & -9 & 6 & 36 & 90 & 198 \\ \hline
\end{tabular}
\end{center}
}{}

\teacheranswerbox{\IconCheck\ ANSWER / ANTICIPATED TALK}{%
Answer key (from TE): No; as the \(x\)-values increase, the \(y\)-values also increase. (This is a direct variation, where \(k = 18\) since \(y/x = 18\) for every pair.)
}

\bankitem{Try It 2 --- Write the equation}{%
Try It 2. In an inverse variation, \(x = 6\) when \(y = 1/2\). (a) What is the equation that represents the inverse variation? (b) What is the value of \(y\) when \(x = 15\)?
}{}

\teacheranswerbox{\IconCheck\ ANSWER / ANTICIPATED TALK}{%
Answer key (confirmed by teacher from TE): (a) \(y = 3/x\) (since \(k = 6\cdot(1/2) = 3\)); (b) when \(x = 15\), \(y = 3/15 = 1/5\).
}

\bankitem{Practice \#13 --- Generalize: write \(k\) in terms of \(x\) and \(y\)}{%
13. Generalize: For an inverse variation, write an equation that gives the value of \(k\) in terms of \(x\) and \(y\).
}{}

\teacheranswerbox{\IconCheck\ ANSWER / ANTICIPATED TALK}{%
Answer key (from TE): \(k = xy\).
}

\bankitem{Practice \#16 --- Apply the equation}{%
16. If \(x\) and \(y\) vary inversely and \(x = 3\) when \(y = 2/3\), what is the value of \(y\) when \(x = -1\)?
}{}

\teacheranswerbox{\IconCheck\ ANSWER / ANTICIPATED TALK}{%
Answer key (from TE): When \(x = -1\), \(y = -2\). IMPORTANT: the original screenshot had the \(y\)-value in the setup cut off at the line break. Reconstructed as \(y = 2/3\) from answer key (\(k = 3*(2/3) = 2\); at \(x = -1\), \(y = -2\) checks). Visual\_needs\_cleanup=true flags that the source screenshot is incomplete.
}

\bankitem{Practice \#10 --- Construct Arguments (extension)}{%
10. Construct Arguments: Explain why zero cannot be in the domain of an inverse variation.
}{}

\teacheranswerbox{\IconCheck\ ANSWER / ANTICIPATED TALK}{%
Answer key (from TE): Division by zero is undefined. (Elaboration: since \(y = k/x\), the value \(x = 0\) produces \(y = k/0\) which is undefined, so \(x = 0\) must be excluded from the domain.)
}

\sectionbanner{SHARE / SUMMARY (5 min)}

\begin{sentenceframebox}
\textbf{Sentence frames}

Pair share: ``As \_\_\_ increases, \_\_\_ decreases. Their product is \_\_\_, so the equation is \(y = \_\_\_/x\).''\\
Whole class: one pair reads their sentence. Class signals thumbs up/sideways.
\end{sentenceframebox}

\phasetag{academic conversation}
\frameworkphaseheader{Share / Summary}{2}{5}{%
\phaseitems{
  \item Cold-call 2 pairs to read their sentence-frame sentence.
  \item Write one student-generated sentence on board.
  \item Preview P2: ``Next class we'll solve real-world inverse-variation problems --- like a bouzouki string.''
}%
}{%
\phaseitems{
  \item Presenting pair reads their sentence.
  \item Class signals thumbs up/sideways.
  \item Note: next class applies this to a real situation.
}%
}{%
\phaseitems{
  \item What's a sentence that captures today's rule?
  \item What's the difference between inverse and direct variation?
}%
}{Cold-call 2 pairs. Keep it brief --- 5 min.}

\phasetag{summary recap (not graded)}
\frameworkphaseheader{Exit Ticket}{1--2}{2}{%
\phaseitems{
  \item Collect at door.
  \item Scan \#2 (inverse vs direct) --- misconceptions tell you what to re-hit in P2 Do Now.
}%
}{%
\phaseitems{
  \item Complete 3 sentence stems.
  \item Be honest about confusion.
}%
}{%
\phaseitems{
  \item Did students correctly distinguish inverse (\(xy = k\)) from direct (\(y/x = k\))?
}%
}{Collect. No grade.}

\begin{callout}{\IconClipboard\ EXIT TICKET --- Student Form}{calloutyellow}
Today I learned that \_\_\_.\\
The difference between inverse and direct variation is \_\_\_.\\
One thing I'm still unsure about is \_\_\_.
\end{callout}

\sectionbanner{OPTIONAL REINFORCEMENT (not collected)}
\textit{If you finish Explore early or want extra practice before Period 2.}

\begin{callout}{\IconSearch\ OPTIONAL \textperiodcentered\ Model \& Discuss Parts B + C}{calloutblue}
Return to the two rectangles from the Do Now (both have area 144). With your partner:\\

B) Use Structure. Considering rectangles with area 144, what happens to the WIDTH as the LENGTH increases?\\

C) Examine at least five more pairs of rectangles with area 144. How would you describe the relationship between their lengths and widths --- in your own words?
\end{callout}

\bankitem{Optional \#1 (Savvas Practice)}{%
15. Do the tables of values represent inverse variations? Explain. See Example 1. Table - \(x: 1, 2, 3, 4, 5, 6;\ y: 60, 30, 20, 15, 12, 10.\)
\begin{center}
\renewcommand{\arraystretch}{1.15}
\begin{tabular}{|c|c|c|c|c|c|c|}
\hline
\(x\) & 1 & 2 & 3 & 4 & 5 & 6 \\ \hline
\(y\) & 60 & 30 & 20 & 15 & 12 & 10 \\ \hline
\end{tabular}
\end{center}
}{}

\teacheranswerbox{\IconCheck\ ANSWER / ANTICIPATED TALK}{%
Answer key (from TE): Yes; as the \(x\)-values increase, the \(y\)-values decrease. Also the product of each pair of \(x\)- and \(y\)-values is 60 (\(k = 60\), so \(y = 60/x\)).
}

\clearpage

\begin{callout}{\IconClock\ PERIOD A \& F --- 65-MIN CLASS NOTES}{calloutpink}
Both sections get 65 min for P1 (Fri F / Mon A). Use the extra 10 min on Explore, NOT Launch.\\
Suggested add: project Savvas Practice \#9 and \#14 on screen (not in packet) for 2 more TPS rounds.\\
Or: extend Model \& Discuss Parts B + C with ``find the equation in the form \(y = k/x\) for your three rectangles.'' (\(k = 144\), so \(y = 144/x\).)
\end{callout}

\begin{callout}{\IconPuzzle\ MODIFICATIONS / SUPPORT FOR IEP STUDENTS}{calloutiep}
\textbullet\ Provide the Inverse Variation Rules card as a sticker/cut-out on their desk.\\
\textbullet\ For Try-It 1, allow students to compute \(xy\) with a calculator for every column.\\
\textbullet\ Accept verbal sentence-frame completions in lieu of written.
\end{callout}

\begin{callout}{\IconGlobe\ SUPPORTS FOR ELL STUDENTS}{calloutell}
\textbullet\ Sentence frames printed on student packet (don't skip).\\
\textbullet\ Word card: ``varies inversely'' = as one goes up, the other goes down proportionally.\\
\textbullet\ ``constant of variation'' = \(k\) = the fixed product \(xy\).\\
\textbullet\ ``product'' = the result of multiplying two numbers.\\
\textbullet\ Pair ELL students with English-dominant partners for TPS.
\end{callout}

\end{document}
